% from ega4-2-fr.tex
\oldpage[IV-2]{128}

est irréductible et de dimension 1, $A/tA$ ne vérifie pas $(S_1)$ et
$\operatorname{Spec}(A)$ n'est pas équidimensionnel.
\end{remark}

\begin{corollary}[5.12.4]
\label{IV.2.s5.12.4-fr}
Soient $A$ un anneau local noethérien caténaire, $t$ un élément régulier
appartenant à l'idéal maximal de $A$, $k$ un entier $\geq1$. Si $A/tA$
vérifie $(S_k)$, il en est de même de $A$.
\end{corollary}

Si $k=1$, le corollaire résulte de (3.4.4), vu l'interprétation de la condition
$(S_1)$ (5.7.5). Si $k\geq2$, il résulte du critère de Hartshorne (5.10.9) que
$A/tA$ est équidimensionnel, et on peut appliquer (5.12.2).

\begin{proposition}[5.12.5]
\label{IV.2.s5.12.5-fr}
Soient $A$ un anneau local noethérien caténaire, $t$ un élément régulier de $A$
appartenant à l'idéal maximal de $A$, $k$ un entier $\geq0$. Si l'anneau
$A/tA$ est réduit, équidimensionnel et vérifie la propriété $(R_k)$, l'anneau
$A$ possède aussi ces trois propriétés.
\end{proposition}

On sait déjà (3.4.6) que $A$ est réduit, et il résulte de (5.12.2), appliqué
pour $k=1$, que $A$ est équidimensionnel. Posons
$X=\operatorname{Spec}(A)$, $X_0=V(t)=\operatorname{Spec}(A/tA)$, et soit $Z$
l'ensemble des points $x\in X$ où $X$ n'est pas régulier; en vertu de (0,
17.3.2) toute spécialisation d'un point de $Z$ appartient à $Z$. Il résulte
d'autre part de (0, 17.1.8) que pour tout point $x\in X_0$ où $X_0$ est
régulier, $X$ est aussi régulier, donc l'ensemble $Z'_0$ des points où $X_0$
n'est pas régulier contient $Z_0=Z\cap X_0$; l'hypothèse entraîne donc que
\[
k<\operatorname{codim}(Z'_0,X_0)\leq\operatorname{codim}(Z_0,X_0).
\]
Mais, puisque $A$ est équidimensionnel et caténaire, on a, d'après (5.12.1),
\[
\operatorname{codim}(Z_0,X_0)\leq\operatorname{codim}(Z,X)
\]
ce qui prouve que $X$ vérifie $(R_k)$.

\begin{remark}[5.12.6]
\label{IV.2.s5.12.6-fr}
Si, dans l'énoncé de (5.12.5), on suppose seulement que $A/tA$ soit réduit et
possède la propriété $(R_k)$, il n'en résulte pas nécessairement que $A$
possède la propriété $(R_k)$. Par exemple, soient $k$ un corps, $P(U,V)$ un
polynôme irréductible de $k[U,V]$ tel que la courbe $\Gamma$ définie par
l'idéal principal $(P)$ dans le plan affine $\operatorname{Spec}(k[U,V])$ ait
un point singulier correspondant à l'idéal maximal $(U)+(V)$ (par exemple
$P(U,V)=U(U^2+V^2)+(U^2-V^2)$, qui donne une cubique à point double). Soient
alors $B$ l'anneau de polynômes à 4 indéterminées $k[T,U,V,W]$, $C$ l'anneau
local de $B$ en l'idéal maximal $BT+BU+BV+BW$, et soient
\[
\mathfrak p=CW+CP(U-T,V),\qquad \mathfrak q=CU.
\]
Considérons alors l'anneau local $A=C/\mathfrak p\mathfrak q$
(géométriquement, si $X$ est le sous-préschéma fermé de l'espace affine à 4
dimensions, formé d'un hyperplan $H$ et d'un « cylindre » $L$ à 2 dimensions
ayant pour « base » la courbe $\Gamma$ et dont la « droite singulière » $Y$
n'est pas contenue dans l'hyperplan $H$, alors $A$ est l'anneau local de $X$ au
point $x$ où $Y$ rencontre $H$). On voit alors aussitôt que
$\operatorname{Spec}(A/tA)$ (où $t$ est l'image canonique de $T$ dans $A$) a
pour seul point singulier $x$, dont l'anneau local est $A/tA$ lui-même, qui est
réduit et de dimension 2; au contraire, le point générique $y$ de la « droite
singulière » $Y$ (définie par l'image dans $A$ de l'idéal
$C(U-T)+CV+CW$) est un point singulier de $X$, et
$\mathcal O_{X,y}$ est de dimension 1; autrement dit, $A/tA$ est réduit et
vérifie $(R_1)$, mais $A$ ne vérifie pas $(R_1)$.
\end{remark}

